\subsubsection{IP Interrupt Status Register (IPISR)}

Address offset: \textbf{0x0C}
\newline
Reset value: 0x0000 0000

\begin{figure}[H]
    \centering
    \begin{bytefield}[bitwidth=\bflenhalfword, endianness=big]{16}
        \bitheader[lsb=16]{16-31} \\
        \bitbox{16}{Reserved}
    \end{bytefield}
\end{figure}

\begin{figure}[H]
    \centering
    \begin{bytefield}[bitwidth=\bflenhalfword, endianness=big]{16}
        \bitheader{0-15} \\
        \bitbox{7}{Reserved} &
        \bitbox{1}{SRA} &
        \bitbox{7}{Reserved} &
        \bitbox{1}{SLA} \\
        \bitbox[]{7}{} &
        \bitbox[]{1}{R} &
        \bitbox[]{7}{} &
        \bitbox[]{1}{R}
    \end{bytefield}
\end{figure}

\begin{itemize}
    \item \textbf{Bit 31:9 - Reserved Bits}
    \item \textbf{Bit 8 - SRA: Sample Right Available}\\
    Dieses Bit wird 1 gelesen, wenn die Daten im ADATRR Register gültig sind.
    Es wird automatisch zurückgesetzt, wenn ein Lesezugriff auf das ADATRR Register ausgeführt wird.
    Dieses Bit löst, wenn das GIE Bit im GIER Register sowie das SRIE Bit im IPIER Register gesetzt ist, einen Interrupt aus.
    \item \textbf{Bit 7:1 - Reserved Bits}
    \item \textbf{Bit 0 - SLA: Sample Left Available}\\
    Dieses Bit wird 1 gelesen, wenn die Daten im ADATLR Register gültig sind.
    Es wird automatisch zurückgesetzt, wenn ein Lesezugriff auf das ADATLR Register ausgeführt wird.
    Dieses Bit löst, wenn das GIE Bit im GIER Register sowie das SLIE Bit im IPIER Register gesetzt ist, einen Interrupt aus.
\end{itemize}